\chapter{第三部分小结}
\label{sec:ch18_summary}

本书的第三部分全景式地探讨了\term{抽象}在系统化程序设计中的核心地位。抽象具有一体两面的特质：\textbf{抽象的创建}与\textbf{抽象的运用}。我们可以将全书第三部分凝聚为两堂最为核心的必修课：

\begin{enumerate}
  \item \textbf{重复的代码模式必须提炼为抽象}。进行抽象意味着提炼出代码中反复出现的共同逻辑，并将所有的具体差异转化为参数。通过精心设计恰当的抽象，程序员为系统构建了稳固的\term{单点控制原则}，从而为未来的演进与维护省去了无数无休止的修补与烦恼。一旦发现逻辑缺陷、性能瓶颈或架构调整，只需在抽象定义这一处进行修改，便可一劳永逸地全盘修正。相反，盲目复制代码意味着埋下定时炸弹，在排查问题时必须大海捞针般搜寻所有副本，极易顾此失彼。
  \item \textbf{大多数现代编程语言都提供了丰富的内置抽象与高阶工具库}。优秀的程序员必须像工匠熟悉刻刀一样熟悉这些核心抽象（\code{map}、\code{filter}、\code{foldr}、\code{build-list} 等）。为了使他人能够高效复用抽象，抽象的作者必须提供详尽严格的设计配方契约——准确的类型签名、清晰的目的说明以及完备的功能用例；而抽象的调用者正是基于这些签名契约，借助局部定义（\code{local}）与匿名函数（\code{lambda}）行云流水地组合出高质量的复杂应用。
\end{enumerate}

除却抽象配方之外，第三部分还注入了一个颠覆传统认知的伟大理念：
\begin{center}
  \textbf{\large\color{NavyTheme}函数本身就是一等公民的值，且函数可以直接承载与表示信息。}
\end{center}

尽管这一思想对于 Lisp 语言家族（如 ISL+）以及现代编程语言理论学家而言由来已久，但在工业界，直到最近十余年，这一强大的范式才终于被绝大多数现代主流编程语言（包括 Python、JavaScript、Java、C\#、C++、Rust 等）彻底拥抱与普及。掌握了函数作为一等公民的高阶抽象思维，你便真正拿到了通往现代高级计算体系的金钥匙。
